\section{DISCUSSION / 讨论}

歪比巴波 歪比巴波 歪比巴波 歪比巴波 歪比巴波 歪比巴波 歪比巴波 歪比巴波 歪比巴波 歪比巴波。

我的刀盾 我的刀盾 我的刀盾 我的刀盾 我的刀盾 我的刀盾 我的刀盾 我的刀盾 我的刀盾 我的刀盾。

咕咕嘎嘎 咕咕嘎嘎 咕咕嘎嘎 咕咕嘎嘎 咕咕嘎嘎 咕咕嘎嘎 咕咕嘎嘎 咕咕嘎嘎 咕咕嘎嘎 咕咕嘎嘎。

巴巴波伊 巴巴波伊 巴巴波伊 巴巴波伊 巴巴波伊 巴巴波伊 巴巴波伊 巴巴波伊 巴巴波伊 巴巴波伊。

\subsection{测试公式与牛顿定理}

为了测试数学公式的排版效果，我们先给出一个经典的“牛顿二项式定理”形式：
\[
(a + b)^n = \sum_{k=0}^{n} \frac{n!}{k!\,(n-k)!}\, a^{n-k} b^{k},
\]
其中 \(n\) 为非负整数，\(\frac{n!}{k!\,(n-k)!}\) 表示组合数。

\subsubsection{心形曲线公式示例}

接下来给出一个常见的“心形公式”，用于绘制爱心形状的曲线：
\[
(x^2 + y^2 - 1)^3 = x^2 y^3,
\]
在适当的坐标范围内，这条隐式曲线会呈现出典型的心形轮廓，非常适合作为浪漫风格的数学测试图像。

\subsection{高雅梗对照表示例}

为了进一步测试表格排版，这里随手给出一张“高雅人士梗”对照表，如表\ref{tab:memes} 所示。

\begin{table}[htbp]
  \centering
  \caption{高雅人士随机玩梗对照表}
  \label{tab:memes}
  \begin{tabular}{ll}
    \hline
    梗 & 备注说明 \\
    \hline
    歪比巴波 & 看到这里说明已经放弃严肃性 \\
    我的刀盾 & 理论上是防御道具，实际上是情绪价值 \\
    咕咕嘎嘎 & 适用于任何不知道说啥的场合 \\
    巴巴波伊 & 当作 QED 使用也完全说得过去 \\
    qing'zu   & 仅限真正的高雅人士使用 \\
    \hline
  \end{tabular}
\end{table}
